% ===========================================================================
%  Chapitre 15 — Statistique à deux variables
%  PE 2026, composante TRAITEMENT DES DONNÉES — série OSE
% ===========================================================================
\bmtchapitre{Statistique à deux variables}
  {Étudier un nuage de points pondérés, déterminer le point moyen, mesurer la
   corrélation linéaire entre deux caractères, ajuster un nuage par une droite
   de régression obtenue par la méthode des moindres carrés, et l'utiliser
   pour une estimation.}
  {Traitement des données \textbullet\ environ 10 heures}

Une entreprise observe deux grandeurs qui évoluent ensemble : le chiffre
d'affaires et la publicité engagée, la production et la surface cultivée, le
nombre d'années et le volume exporté. La question n'est plus, comme en
première, de résumer une seule série de nombres, mais de comprendre le lien
entre deux séries observées sur les mêmes individus.

Ce chapitre répond à trois questions, toujours dans le même ordre. Les deux
grandeurs sont-elles réellement liées, et avec quelle intensité ? Si le lien
est linéaire, quelle droite le représente le mieux ? Et cette droite
permet-elle de prévoir raisonnablement une valeur non observée ? C'est
exactement le raisonnement que suit un économiste pour construire une
prévision, et c'est celui que demande chaque année l'exercice de statistique
de l'épreuve.

% ---------------------------------------------------------------------------
\begin{revision}
Statistique à une variable de première.

\begin{enumerate}[label=\textbf{\arabic*.}, leftmargin=*, itemsep=0.35em]
  \item Rappeler la formule de la moyenne d'une série $\left(x_{i}\right)$ de
        même effectif.
  \item Rappeler la formule de la variance
        $V = \dfrac{1}{n}\sum x_{i}^{2}-\overline{x}^{2}$.
  \item Que représente géométriquement un nuage de points ?
  \item Rappeler l'équation réduite d'une droite non verticale.
  \item Résoudre le système $\begin{cases} 2a+b = 7 \\ 5a+b = 13
        \end{cases}$.
\end{enumerate}
\end{revision}

\begin{automatismes}
Sans calculatrice, sans brouillon.

\begin{enumerate}[label=\textbf{\alph*)}, leftmargin=*, itemsep=0.25em]
  \item Moyenne de $1,2,3,4,5$
  \item Moyenne de $10,20,30$
  \item Coefficient directeur de la droite passant par $(0\,;1)$ et $(2\,;7)$
  \item Ordonnée à l'origine de $y = 3x+2$ en $x = 0$
  \item Si $r = 0{,}97$, la corrélation est-elle forte ou faible ?
  \item Si $r = 0{,}1$, la corrélation est-elle forte ou faible ?
  \item Signe de $r$ si le nuage monte de gauche à droite
  \item $\overline{x}$ pour $x_{i} \in \{1,2,3,4,5,6\}$
\end{enumerate}
\end{automatismes}

\begin{corrige}[automatismes]
a) $3$ \quad b) $20$ \quad c) $3$ \quad d) $2$ \quad e) forte \quad
f) faible \quad g) positif \quad h) $3{,}5$.
\end{corrige}

% ===========================================================================
\section{Série statistique double et nuage de points}

\begin{definition}[série statistique à deux variables]
Une \textbf{série statistique double} associe à $n$ individus deux
caractères quantitatifs, mesurés par les couples
$\left(x_{1}\,;y_{1}\right),\dots,\left(x_{n}\,;y_{n}\right)$. Le
\textbf{nuage de points} est l'ensemble des points $M_{i}\left(x_{i}\,;
y_{i}\right)$ dans un repère du plan.
\end{definition}

\begin{definition}[point moyen]
Le \textbf{point moyen} du nuage est le point $G\left(\overline{x}\,;
\overline{y}\right)$ où
\[
  \overline{x} = \frac{1}{n}\sum_{i=1}^{n}x_{i},
  \qquad
  \overline{y} = \frac{1}{n}\sum_{i=1}^{n}y_{i} .
\]
\end{definition}

\begin{remarque}
Lorsque le nuage prend l'allure d'une bande allongée autour d'une droite, on
dit qu'un \textbf{ajustement affine} est envisageable : c'est ce que
suggère l'œil, et c'est ce que les outils de ce chapitre vont confirmer ou
contredire par le calcul.
\end{remarque}

% ===========================================================================
\section{Covariance et coefficient de corrélation linéaire}

\begin{definition}[covariance]
La \textbf{covariance} des séries $\left(x_{i}\right)$ et
$\left(y_{i}\right)$ est
\[
  \Cov(x,y) = \frac{1}{n}\sum_{i=1}^{n}x_{i}y_{i}-\overline{x}\,\overline{y} .
\]
\end{definition}

\begin{definition}[coefficient de corrélation linéaire]
Le \textbf{coefficient de corrélation linéaire} est
\[
  r = \frac{\Cov(x,y)}{\sigma_{x}\,\sigma_{y}},
\]
où $\sigma_{x}$ et $\sigma_{y}$ sont les écarts-types des deux séries. Il
vérifie toujours $-1 \leq r \leq 1$.
\end{definition}

\begin{propriete}[lecture de $r$]
\begin{itemize}[leftmargin=*, itemsep=0.1em]
  \item $\abs r$ proche de $1$ : les points sont presque alignés, la liaison
        linéaire est forte.
  \item $\abs r$ proche de $0$ : le nuage n'a pas de forme allongée, ou sa
        forme n'est pas linéaire.
  \item Le signe de $r$ indique le sens de variation : positif si $y$
        augmente en moyenne avec $x$, négatif sinon.
\end{itemize}
En pratique, un ajustement affine n'est retenu que si $\abs r$ dépasse
environ $0{,}87$ ; l'énoncé précise parfois le seuil à partir duquel il
autorise l'ajustement.
\end{propriete}

\begin{avertissement}
Une forte corrélation ne prouve pas un lien de cause à effet. Deux séries
peuvent varier ensemble sans que l'une explique l'autre — un troisième
facteur, ou le simple hasard sur peu d'observations, peut produire un $r$
élevé. Le coefficient mesure une liaison statistique, pas une causalité
économique.
\end{avertissement}

\begin{exemple}[calcul complet de $r$]
Une pépinière relève, sur cinq campagnes, le nombre $x_{i}$ de pieds
plantés (en centaines) et la masse $y_{i}$ récoltée (en tonnes) :
\[
  \begin{array}{c|ccccc}
    x_{i} & 1 & 2 & 3 & 4 & 5 \\
    \hline
    y_{i} & 3 & 5 & 6 & 8 & 9
  \end{array}
\]
On a $\overline{x} = 3$, $\overline{y} = 6{,}2$,
$\dfrac{1}{5}\sum x_{i}y_{i}
= \dfrac{3+10+18+32+45}{5} = \dfrac{108}{5} = 21{,}6$, d'où
\[
  \Cov(x,y) = 21{,}6-3 \times 6{,}2 = 3.
\]
Par ailleurs $\dfrac{1}{5}\sum x_{i}^{2} = \dfrac{55}{5} = 11$, donc
$\sigma_{x}^{2} = 11-9 = 2$ ; et $\dfrac{1}{5}\sum y_{i}^{2}
= \dfrac{9+25+36+64+81}{5} = \dfrac{215}{5} = 43$, donc
$\sigma_{y}^{2} = 43-6{,}2^{2} = 4{,}56$. Finalement
\[
  r = \frac{3}{\sqrt2 \times \sqrt{4{,}56}} \approx \frac{3}{3{,}019}
  \approx 0{,}994 .
\]
La corrélation est très forte : un ajustement affine est pleinement
justifié.
\end{exemple}

% ===========================================================================
\section{Ajustement affine par la méthode des moindres carrés}

\begin{theoreme}[droite de régression de $y$ en $x$]
Lorsque $\abs r$ est proche de $1$, la droite $(D)$ qui approche le mieux le
nuage — au sens où elle minimise la somme des carrés des écarts verticaux —
a pour équation
\[
  y = ax+b, \qquad
  a = \frac{\Cov(x,y)}{\sigma_{x}^{2}}, \qquad
  b = \overline{y}-a\,\overline{x} .
\]
Cette droite passe toujours par le point moyen $G$.
\end{theoreme}

\begin{methode}{déterminer la droite de régression de $y$ en $x$}
\begin{enumerate}[label=\textbf{\arabic*.}, leftmargin=*, itemsep=0.15em]
  \item Calculer $\overline{x}$, $\overline{y}$, puis
        $\dfrac{1}{n}\sum x_{i}y_{i}$ et $\dfrac{1}{n}\sum x_{i}^{2}$.
  \item En déduire $\Cov(x,y)$ et $\sigma_{x}^{2}$.
  \item Calculer $a = \dfrac{\Cov(x,y)}{\sigma_{x}^{2}}$, puis
        $b = \overline{y}-a\,\overline{x}$.
  \item Pour estimer $y$ à partir d'une valeur de $x$ non observée, on
        remplace simplement $x$ dans l'équation de $(D)$ — en gardant à
        l'esprit qu'une estimation n'est fiable que proche du domaine
        observé.
\end{enumerate}
\end{methode}

\begin{visualisation}[nuage, point moyen et droite de régression]
\begin{center}
\begin{tikzpicture}[x=1.55cm, y=0.42cm, font=\footnotesize]
  \draw[bmtGris, -{Stealth[length=2.6mm]}] (-0.4,0) -- (6.3,0)
    node[anchor=north west] {$x$ (rang)};
  \draw[bmtGris, -{Stealth[length=2.6mm]}] (0,-0.5) -- (0,11)
    node[anchor=south east] {$y$};
  \foreach \x/\y in {1/3, 2/5, 3/6, 4/8, 5/9}
    \filldraw[bmtIndigo] (\x,\y) circle (2.1pt);
  \filldraw[bmtLaterite] (3,6.2) circle (2.6pt);
  \node[bmtLaterite, anchor=west] at (3.15,6.2) {$G$};
  \draw[bmtLaterite, line width=1pt]
    (0.3,3.1) -- (5.7,9.7);
  \node[bmtGris, anchor=west] at (5.75,9.6) {$(D)$};
\end{tikzpicture}
\end{center}

Le nuage s'étire visiblement le long d'une direction : c'est le signe qu'un
ajustement affine a du sens. La droite $(D)$ ne cherche pas à passer par tous
les points — c'est impossible en général — mais à minimiser globalement
l'écart vertical entre chaque point et la droite. Elle passe, elle,
exactement par le point moyen $G$ : c'est la propriété à retenir pour vérifier
un calcul.
\end{visualisation}

\begin{remarque}[méthode de Mayer]
Une méthode plus rapide, mais plus approximative, consiste à partager le
nuage en deux groupes de tailles égales selon les $x_{i}$ croissants, à
calculer le point moyen $G_{1}$ du premier groupe et $G_{2}$ du second, puis
à prendre pour droite d'ajustement la droite $\left(G_{1}G_{2}\right)$. Cette
\textbf{méthode de Mayer} donne une droite proche de la droite des moindres
carrés lorsque le nuage est bien aligné, mais elle n'a pas la même
justification théorique et reste moins précise.
\end{remarque}

\begin{entrainement}[régression et corrélation]
\exo[\niveau{1}] Une série de cinq points a pour point moyen
$G(4\,;10)$ et coefficient directeur de régression $a = 1{,}5$. Donner
l'équation de $(D)$.

\exo[\niveau{2}] Une coopérative relève sur quatre ans :
\[
  \begin{array}{c|cccc}
    x_{i} \text{ (rang)} & 1 & 2 & 3 & 4 \\
    \hline
    y_{i} \text{ (tonnes)} & 20 & 24 & 27 & 33
  \end{array}
\]
Calculer $\overline{x}$, $\overline{y}$, $\Cov(x,y)$ et $\sigma_{x}^{2}$,
puis l'équation de la droite de régression de $y$ en $x$.

\exo[\niveau{2}] Avec les données précédentes, estimer la production
attendue au rang $x = 6$. Cette estimation vous paraît-elle également
fiable qu'une estimation au rang $x = 5$ ? Justifier.

\exo[\niveau{3}] Montrer, à partir des formules du cours, que la droite de
régression de $y$ en $x$ passe toujours par le point moyen $G$.
\end{entrainement}

% ===========================================================================
\begin{annales}
La statistique à deux variables clôt souvent l'exercice 2 de l'épreuve, sur
un scénario stable : un tableau de données économiques sur plusieurs années,
un calcul de corrélation, une droite de régression, puis une estimation.
L'énoncé qui suit est repris d'un sujet réellement posé.

\exoann{Baccalauréat de l'Enseignement Général, Madagascar, série OSE, option OSE, session 2026 — exercice 2}
Le tableau suivant donne la production $y_{i}$ (en tonnes) de maïs d'un Petit
Périmètre Irrigué, observée de 2019 à 2024, $x_{i}$ désignant le rang de
l'année :
\[
  \begin{array}{c|cccccc}
    x_{i} & 1 & 2 & 3 & 4 & 5 & 6 \\
    \hline
    y_{i} & 60 & 65 & 70 & 74 & 75 & 80
  \end{array}
\]
On admet que $\sum x_{i} = 21$, $\sum y_{i} = 424$, $\sum x_{i}^{2} = 91$,
$\sum y_{i}^{2} = 30\,226$ et $\sum x_{i}y_{i} = 1\,551$.
\begin{enumerate}[label=\textbf{\arabic*.}, leftmargin=*, itemsep=0.2em]
  \item Représenter le nuage de points et placer le point moyen $G$.
  \item Calculer le coefficient de corrélation linéaire de la série, et
        interpréter le résultat.
  \item Déterminer, par la méthode des moindres carrés, une équation de la
        droite de régression $(D)$ de $y$ en $x$.
  \item En supposant que l'évolution se poursuive selon le même modèle,
        estimer la production attendue en 2028.
\end{enumerate}

\exoann[ajustement affine et ajustement logarithmique]{Baccalauréat, France, série ES, Pondichéry, 17 avril 2012 — exercice 2}
Anna a créé un site web. Le tableau ci-dessous donne le nombre hebdomadaire
de visiteurs au cours des huit premières semaines suivant sa création, $x_{i}$
désignant le rang de la semaine :
\[
  \begin{array}{c|cccccccc}
    x_{i} & 1 & 2 & 3 & 4 & 5 & 6 & 7 & 8 \\
    \hline
    y_{i} & 205 & 252 & 327 & 349 & 412 & 423 & 441 & 472
  \end{array}
\]
\begin{enumerate}[label=\textbf{\arabic*.}, leftmargin=*, itemsep=0.2em]
  \item Représenter le nuage de points $M_{i}(x_{i}\,;y_{i})$, puis
        déterminer les coordonnées du point moyen $G$ (l'ordonnée sera
        arrondie à l'unité) et le placer sur le graphique.
  \item Déterminer, par la méthode des moindres carrés, une équation
        $y = ax+b$ de la droite d'ajustement affine de $y$ en $x$ ($a$ et $b$
        arrondis à l'entier le plus proche), puis l'utiliser pour estimer le
        nombre de visiteurs lors de la dixième semaine.
  \item L'augmentation du nombre de visiteurs étant plus faible sur les
        dernières semaines, on envisage un ajustement de type logarithmique
        en posant $z = \ln(x)$.
        \begin{enumerate}[label=\textbf{\alph*)}, leftmargin=1.2em]
          \item Calculer les valeurs manquantes $z_{3}$ et $z_{6}$
                (arrondies à $10^{-3}$) dans le tableau
                $z_{i} = \ln(x_{i})$.
          \item On admet que la droite d'ajustement affine de $y$ en $z$,
                obtenue par la méthode des moindres carrés, a pour équation
                $y = 133z+184$. Donner, à l'aide de cet ajustement, une
                nouvelle estimation (arrondie à l'unité) du nombre de
                visiteurs lors de la dixième semaine.
          \item Déterminer, à l'aide de ce second ajustement, le rang de la
                semaine à partir duquel le nombre prévisible de visiteurs
                dépasse $600$.
        \end{enumerate}
\end{enumerate}
\end{annales}

\begin{corrige}[annales]
\textbf{1--2.} $\overline{x} = \dfrac{21}{6} = 3{,}5$ et
$\overline{y} = \dfrac{424}{6} \approx 70{,}67$ : le point moyen est
$G(3{,}5\,;70{,}67)$.

En posant $S_{xy} = \sum x_{i}y_{i}-\dfrac{\left(\sum x_{i}\right)
\left(\sum y_{i}\right)}{n}$, $S_{xx} = \sum x_{i}^{2}
-\dfrac{\left(\sum x_{i}\right)^{2}}{n}$ et $S_{yy} = \sum y_{i}^{2}
-\dfrac{\left(\sum y_{i}\right)^{2}}{n}$ :
\[
  S_{xy} = 1\,551-\frac{21 \times 424}{6} = 1\,551-1\,484 = 67,
\]
\[
  S_{xx} = 91-\frac{441}{6} = 17{,}5, \qquad
  S_{yy} = 30\,226-\frac{179\,776}{6} \approx 263{,}33 .
\]
D'où
\[
  r = \frac{S_{xy}}{\sqrt{S_{xx}\,S_{yy}}}
  = \frac{67}{\sqrt{17{,}5 \times 263{,}33}}
  \approx \frac{67}{67{,}88} \approx 0{,}987 .
\]
Ce coefficient, très proche de $1$, indique une corrélation linéaire très
forte et positive : la production croît régulièrement d'année en année, et un
ajustement affine est pleinement justifié.

\textbf{3.} Le coefficient directeur vaut
$a = \dfrac{S_{xy}}{S_{xx}} = \dfrac{67}{17{,}5} \approx 3{,}83$, et
l'ordonnée à l'origine
$b = \overline{y}-a\,\overline{x} \approx 70{,}67-3{,}83 \times 3{,}5
\approx 57{,}27$. Une équation de $(D)$ est donc
\[
  y \approx 3{,}83\,x+57{,}27 .
\]

\textbf{4.} L'année 2028 correspond au rang $x = 10$ (2019 étant le rang
$1$). En reportant dans l'équation de $(D)$ :
\[
  y \approx 3{,}83 \times 10+57{,}27 \approx 95{,}6 \text{ tonnes}.
\]
La coopérative peut donc anticiper une production voisine de $95{,}6$
tonnes en 2028 — étant entendu que l'estimation devient moins fiable à
mesure que l'on s'éloigne des années réellement observées.

\textbf{5.} $\overline{x} = \dfrac{36}{8} = 4{,}5$ et la somme des $y_{i}$
vaut $2\,881$, donc $\overline{y} = \dfrac{2\,881}{8} \approx 360$ : le point
moyen est $G(4{,}5\,;360)$.

\textbf{6.} Avec $\sum x_{i} = 36$, $\sum y_{i} = 2\,881$,
$\sum x_{i}^{2} = 204$ et $\sum x_{i}y_{i} = 14\,547$,
\[
  a = \frac{8\times14\,547-36\times2\,881}{8\times204-36^{2}}
  = \frac{12\,660}{336} \approx 37{,}7 \approx 38,
\]
et $b = \overline{y}-a\,\overline{x} \approx 360{,}125-37{,}7\times4{,}5
\approx 190{,}6 \approx 191$. La droite d'ajustement est donc
$y \approx 38x+191$, et pour $x = 10$ :
$y \approx 38\times10+191 = 571$ visiteurs.

\textbf{7. a)} $z_{3} = \ln3 \approx 1{,}099$ et $z_{6} = \ln6 \approx
1{,}792$.
\textbf{b)} Pour $x = 10$, $z = \ln10 \approx 2{,}303$, et
$y \approx 133\times2{,}303+184 \approx 490$ visiteurs.
\textbf{c)} L'inéquation $133z+184>600$ équivaut à
$z > \dfrac{416}{133} \approx 3{,}128$, soit
$\ln(x) > 3{,}128$, c'est-à-dire $x > e^{3{,}128} \approx 22{,}8$. Le nombre
prévisible de visiteurs dépasse donc $600$ à partir de la $23^{\text{e}}$
semaine.
\end{corrige}

% ===========================================================================
\begin{jecherche}[les exportations de la coopérative caféière]
Une coopérative de producteurs de café relève, sur six campagnes
consécutives, la quantité $y_{i}$ exportée (en tonnes) :
\[
  \begin{array}{c|cccccc}
    x_{i} \text{ (rang)} & 1 & 2 & 3 & 4 & 5 & 6 \\
    \hline
    y_{i} \text{ (tonnes)} & 50 & 54 & 60 & 63 & 69 & 74
  \end{array}
\]

\begin{enumerate}[label=\textbf{\arabic*.}, leftmargin=*, itemsep=0.3em]
  \item Calculer $\overline{x}$ et $\overline{y}$, puis placer le point
        moyen $G$ sur un nuage de points.
  \item Calculer $\Cov(x,y)$, $\sigma_{x}^{2}$ et le coefficient de
        corrélation linéaire $r$. Un ajustement affine est-il justifié ?
  \item Déterminer l'équation de la droite de régression $(D)$ de $y$ en
        $x$, par la méthode des moindres carrés.
  \item Estimer les exportations attendues à la neuvième campagne.
  \item Un gel tardif détruit une partie de la récolte de la septième
        campagne, ramenant les exportations à $40$ tonnes cette année-là.
        Sans refaire le calcul complet, expliquer qualitativement l'effet de
        cette valeur sur le coefficient de corrélation si elle était
        intégrée à la série.
  \item La coopérative souhaite atteindre $100$ tonnes exportées. À partir de
        quel rang le modèle le prévoit-il ?
\end{enumerate}
\end{jecherche}

\begin{corrige}[les exportations de la coopérative caféière]
\textbf{1.} $\overline{x} = \dfrac{1+2+3+4+5+6}{6} = 3{,}5$ et
$\overline{y} = \dfrac{50+54+60+63+69+74}{6} = \dfrac{370}{6}
\approx 61{,}67$. Le point moyen est $G(3{,}5\,;61{,}67)$.

\textbf{2.} $\sum x_{i}y_{i} = 50+108+180+252+345+444 = 1\,379$, donc
\[
  S_{xy} = 1\,379-\frac{21 \times 370}{6} = 1\,379-1\,295 = 84 .
\]
Comme au chapitre, $S_{xx} = 91-\dfrac{441}{6} = 17{,}5$, et
$\sum y_{i}^{2} = 2\,500+2\,916+3\,600+3\,969+4\,761+5\,476 = 23\,222$, donc
\[
  S_{yy} = 23\,222-\frac{370^{2}}{6} \approx 405{,}33 .
\]
D'où
$r = \dfrac{84}{\sqrt{17{,}5 \times 405{,}33}}
\approx \dfrac{84}{84{,}22} \approx 0{,}997$ : la corrélation est quasiment
parfaite, l'ajustement affine est pleinement justifié.

\textbf{3.} $a = \dfrac{84}{17{,}5} = 4{,}8$ et
$b = 61{,}67-4{,}8 \times 3{,}5 \approx 44{,}87$. Une équation de $(D)$ est
$y \approx 4{,}8\,x+44{,}87$.

\textbf{4.} Au rang $x = 9$ :
$y \approx 4{,}8 \times 9+44{,}87 \approx 88{,}1$ tonnes.

\textbf{5.} Une valeur de $y$ nettement inférieure à ce que prévoit la
tendance générale — ici un accident climatique, non représentatif de
l'évolution normale — éloigne ce point de la droite : la dispersion autour de
l'ajustement augmente, et le coefficient $\abs r$ diminuerait si cette
campagne était intégrée à la série sans précaution particulière.

\textbf{6.} On résout $4{,}8\,x+44{,}87 \geq 100$, soit
$x \geq \dfrac{55{,}13}{4{,}8} \approx 11{,}5$ : le modèle prévoit d'atteindre
$100$ tonnes à partir de la douzième campagne.
\end{corrige}

% ===========================================================================
\begin{jemeteste}
Répondre par vrai ou faux, en justifiant en une ligne.

\begin{enumerate}[label=\textbf{\arabic*.}, leftmargin=*, itemsep=0.28em]
  \item Le point moyen appartient toujours à la droite de régression.
  \item Un coefficient de corrélation négatif signale une erreur de calcul.
  \item $r$ est toujours compris entre $-1$ et $1$.
  \item Une forte corrélation prouve un lien de cause à effet.
  \item Extrapoler loin du domaine observé est toujours fiable, dès lors que
        $r$ est proche de $1$.
  \item La méthode des moindres carrés minimise la somme des écarts
        verticaux au carré entre les points et la droite.
\end{enumerate}
\end{jemeteste}

\begin{corrige}[je me teste]
\textbf{1.} Vrai — c'est une propriété systématique de la droite de
régression.
\textbf{2.} Faux — un coefficient négatif signale simplement une liaison
décroissante.
\textbf{3.} Vrai, par construction du coefficient de corrélation.
\textbf{4.} Faux — la corrélation est une liaison statistique, pas une
preuve de causalité.
\textbf{5.} Faux — l'extrapolation devient de moins en moins fiable à mesure
que l'on s'éloigne des valeurs observées.
\textbf{6.} Vrai, c'est la définition même de l'ajustement par moindres
carrés.
\end{corrige}

% ===========================================================================
\begin{commeaubac}
\bmtsource{Exercice original, au format de l'épreuve — série OSE}
Une entreprise de transformation de litchis relève, sur cinq campagnes
d'exportation consécutives, la masse $y_{i}$ exportée (en tonnes) :
\[
  \begin{array}{c|ccccc}
    x_{i} \text{ (rang)} & 1 & 2 & 3 & 4 & 5 \\
    \hline
    y_{i} \text{ (tonnes)} & 120 & 135 & 150 & 158 & 172
  \end{array}
\]
On donne $\sum x_{i} = 15$, $\sum y_{i} = 735$, $\sum x_{i}^{2} = 55$,
$\sum y_{i}^{2} = 109\,673$.

\begin{enumerate}[label=\textbf{\arabic*.}, leftmargin=*, itemsep=0.25em]
  \item Calculer les coordonnées du point moyen $G$. \hfill\textup{(0,5 pt)}
  \item Calculer le coefficient de corrélation linéaire de la série.
        \hfill\textup{(0,75 pt)}
  \item Déterminer une équation de la droite de régression de $y$ en $x$
        par la méthode des moindres carrés. \hfill\textup{(0,75 pt)}
  \item Estimer la masse exportée à la huitième campagne.
        \hfill\textup{(0,5 pt)}
\end{enumerate}
\end{commeaubac}

\begin{corrige}[comme au baccalauréat]
\textbf{1.} $\overline{x} = \dfrac{15}{5} = 3$ et
$\overline{y} = \dfrac{735}{5} = 147$ : $G(3\,;147)$.

\textbf{2.} $\sum x_{i}y_{i} = 120+270+450+632+860 = 2\,332$, d'où
\[
  S_{xy} = 2\,332-\frac{15 \times 735}{5} = 2\,332-2\,205 = 127,
\]
\[
  S_{xx} = 55-\frac{225}{5} = 10, \qquad
  \sum y_{i}^{2} = 14\,400+18\,225+22\,500+24\,964+29\,584 = 109\,673,
\]
\[
  S_{yy} = 109\,673-\frac{735^{2}}{5} = 109\,673-108\,045 = 1\,628 .
\]
\[
  r = \frac{127}{\sqrt{10 \times 1\,628}} = \frac{127}{\sqrt{16\,280}}
  \approx \frac{127}{127{,}59} \approx 0{,}995 .
\]
La corrélation est quasiment parfaite.

\textbf{3.} $a = \dfrac{127}{10} = 12{,}7$ et
$b = 147-12{,}7 \times 3 = 147-38{,}1 = 108{,}9$. Une équation de $(D)$ est
\[
  y = 12{,}7\,x+108{,}9 .
\]

\textbf{4.} À la huitième campagne, $x = 8$ :
$y = 12{,}7 \times 8+108{,}9 = 101{,}6+108{,}9 = 210{,}5$ tonnes.
\end{corrige}

% ===========================================================================
\begin{notepourlaclasse}
Les élèves confondent souvent les deux formules de la variance et de la
covariance — moyenne des produits moins produit des moyennes. Faites
toujours écrire les deux sommes intermédiaires,
$\frac1n\sum x_{i}y_{i}$ et $\frac1n\sum x_{i}^{2}$, avant tout calcul de
$\Cov$ ou de $\sigma_{x}^{2}$ : c'est l'étape où se produit l'essentiel des
erreurs.

Insistez sur la propriété du point moyen : une droite de régression qui ne
passe pas par $G$ trahit une erreur de calcul quelque part, et c'est le test
le plus rapide pour vérifier un résultat en fin d'exercice.

Enfin, la dernière question — une estimation par extrapolation — mérite
toujours un commentaire qualitatif, pas seulement un nombre. Les meilleurs
sujets, comme celui de 2026, valorisent explicitement la prudence face à une
projection éloignée des données observées : c'est un réflexe économique
autant que mathématique, à faire prendre aux élèves dès ce chapitre.
\end{notepourlaclasse}
